\scanpage{84}
\[
\tag{37}\diagram{\begin{tikzcd}[ampersand replacement=\&]\Pi_1(X)\arrow[r,"f=\Pi_1(f)"]\&\Pi_1(Y)\end{tikzcd}}
\]
\[
\tag{38}\diagram{\begin{tikzcd}[ampersand replacement=\&]\widehat\Pi_1(X)\arrow[r,"\widehat f=\widehat\Pi_1(f)"]\&\widehat\Pi_1(Y)\end{tikzcd}},
\]
dont le deuxième peut être déduit du premier par passage au groupoïde profini associé, mais peut aussi se décrire directement en termes purement schématiques, en termes \inserted{des morphismes} de schémas \deleted{\unsure{sur|au}} $X\mathbin{-}Y$ sur le corps alg. clos $\mathbb C$. Par suite, (38) commute à l'action de $\unsure{\mathbf{\Gamma}|\Gamma}$,
\[
\tag{39}\widehat f({}^ul)={}^u(\widehat f(l))
\]
et si on se borne aux sous-groupoïdes induits sur $X_{\bar K},Y_{\bar K}$, \inserted{\unsure{d'où} ($\widehat\Pi_1\,Y_{\bar K}$)}, ce dernier commute à l'action de $\Gamma=\operatorname{Gal}\bar K/K$. Si on se $=$ aux sous-groupoïdes induits $X_K$, d'où des $f_l(u),g_l(u)$ pour un chemin $l$ dans $X$ ou dans $Y$, on trouve aussitôt
\[
\tag{40}\left\{\begin{aligned}
\widehat f(f_l(u))&=f_{\widehat f(l)}(u)\\
\widehat f(g_l(u))&=g_{\widehat f(l)}(u).
\end{aligned}\right.
\]
\margin{Même formule \unsure{dans le cas} d'un $T$ $1$-connexe sur $X$ ($T$ esp. \unsure{topologique}), \unsure{donc} au-dessus de $Y$. On aura:
\[
\tag{40bis}g^{T/Y}_{y,x}(u)=\widehat f(g^{T/X}_{y,x}(u))
\]
\deleted{\unsure{d'où}}}
\margin{83}
